\section{The \texttt{SANS\_DebyeS} McStas Component}
Sample for Small Angle Neutron Scattering, Debye-Scherrer Ring

\subsection*{Identification}
\begin{itemize}
  \item \textbf{Author:} Henrich Frielinghaus
  \item \textbf{Origin:} FZ-Juelich/FRJ-2/IFF/KWS-2
  \item \textbf{Date:} Sept 2004
\end{itemize}

\subsection*{Description}
Debye-Scherrer Ring: Model for highly ordered diblock copolymer. This example is highly simple. Multiple scattering neglected, but could be incorporated. Sample that only produces a DebyeSherrer ring. Advantages: Transmitted beam is simulated. Intensities still in flux units. Minor point: Absolute scattering probability given by transmission.

Sample components leave the units of flux for the probability of the individual paths. That is more consitent than the Sans\_spheres routine. Furthermore one can simulate the transmitted beam. This allows to determine the needed size of the beam stop. Only absorption has not been included yet to these sample-components. But thats really nothing.

Example: SANS\_DebyeS(transm=0.5, qDS=0.008, xwidth=0.01, yheight=0.01, zdepth=0.001)

\subsection*{Input parameters}
Parameters in \textbf{boldface} are required; the others are optional.

\begin{longtable}{p{0.22\textwidth}p{0.12\textwidth}p{0.46\textwidth}p{0.14\textwidth}}
\toprule
\textbf{Name} & \textbf{Unit} & \textbf{Description} & \textbf{Default} \\
\midrule
\endhead
transm & 1 & coherent transmission of sample for the optical path "zdepth" & 0.5 \\
qDS & \AA{}$^{-1}$ & Scattering vector of Debye-Sherrer ring & 0.008 \\
xwidth & m & horiz. dimension of sample, as a width (m) & 0.01 \\
yheight & m & vert.. dimension of sample, as a height (m) & 0.01 \\
zdepth & m & thickness of sample (m) & 0.001 \\
\bottomrule
\end{longtable}

\subsection*{Links}
\begin{itemize}
  \item Component source code found in file \texttt{SANS\_DebyeS.comp}.
  \item Sans\_spheres component
\end{itemize}
\IfFileExists{contrib/SANS_DebyeS_static.tex}{\input{contrib/SANS_DebyeS_static.tex}}{}